% As written in the PDF: Keep in mind that this is just a template to save you time!
% You can still make it appear as if you like it the most!
\ifLongThesisStyle
\chapter*{Abstract}
\else%
\par\bigskip\textbf{Abstract.}
\fi
%
Non-terminating programs such as streams are ubiquitous in modern software, but unexpected computational demands can make these programs expensive to run because they operate continually. Type systems like graded modal types provide static guarantees about the (amortised) computational cost of programs.  However, such guarantees are typically given for programs that operate on finite data structures. In addition, productivity is a desired property for infinite programs, ensuring that outputs are produced continuously. The later modality allows types to ensure productivity through the idea of guarded recursion, yet it cannot control their costs.

This thesis presents an \emph{integrated} type system that contains a graded cost monad and the later modality. Within this system, types can specify the upper-bound cost for computations over infinite data structures. Moreover, we explore the interaction between the two modalities, which enables us to model costs as functions over an abstract time clock associated with the later modality, allowing programs to express costs that vary from step to step, accounting for differences in the cost of producing successive outputs. We use logical relations to interpret the type system and prove it sound.

Amortised analysis is a common technique that can yield tighter cost bounds than traditional worst-case analysis. To support it, we extend our type system with affine types and a graded potential modality, based on existing work. We restrict potential grades to be constant over the time clock and extend the logical relation to establish soundness for the resulting type system.